\chapter{Implementation}
\label{chap:implementation}

\section{Repository Organisation}

The implementation is organised as experiment packages within a single repository. Each package retains its source applications, Cooja configurations, run logs, validation outputs, extracted features, analysis results, and notes. This structure separates baseline evidence from later diagnostic or robustness work. In particular, the original blackhole and sinkhole evidence is preserved, while additional attacks, Sybil, robustness, drift monitoring, and defence experiments use separate directories.

The Contiki-NG source tree contains small RPL-lite hooks required for behaviours that cannot be expressed solely in an application. Hooks use weak shared variables so that standard applications leave them disabled. An attack application provides a strong variable definition and changes the flag after the activation delay. This pattern limits the change surface and allows attack and control applications to share the same RPL implementation.

\section{Shared Delayed-Activation Pattern}

Each attack router logs that it has started with malicious behaviour disabled. A timer waits for 240 seconds, after which the shared flag is enabled and an activation message is recorded. The process then continues normally rather than terminating. The corresponding control application permanently defines the flag as zero. This pattern makes the activation boundary explicit and supports matched attack-control comparisons.

The pattern also separates the existence of attack code from observed malicious behaviour. Before 240 seconds, attack runs act as an internal baseline. After activation, the hook or application changes only its assigned mechanism. Validation scripts check both the startup and activation phases. The explicit log messages are useful for debugging and evidence validation but are excluded from learning features.

\section{Forwarding and Rank Attacks}

The blackhole implementation drops forwarded application traffic after activation. Its validated post-activation responses fall to approximately zero to three in attack runs. Grayhole applies selective dropping rather than complete suppression, producing 208 to 211 validated drop events per run. These attacks principally affect data-plane delivery and are therefore visible through application and forwarding-related features.

The sinkhole hook changes the rank written into outgoing DIO advertisements to the minimum hop rank increment while preserving the node's internal rank. Preserving internal state is important because it models deceptive advertisement rather than forcing the malicious node to believe its own false claim. The hook logs each advertised-rank manipulation. The increase-rank hook remains a separate conditional branch and advertises degraded rank information. Keeping these branches separate prevents one experiment from silently changing another attack.

\section{Control-Plane and Route-Selection Attacks}

DIS flooding repeatedly sends solicitation messages after activation, producing 149 validated effect events per attack run. DIO suppression changes the availability of routing advertisements. Worst-parent changes parent-selection logic so that an undesirable candidate is preferred, generating thousands of selection events across runs. These attacks exercise control-plane rate, availability, and route-choice features rather than the direct packet-loss signature learned from blackhole.

The wormhole experiment requires a different topology treatment. Two endpoint motes are placed in a directed-radio configuration and a tunnel-like radio path between nodes 16 and 17 is activated at 240 seconds. The final configurations use 17 motes and preserve this experiment separately from the forced-path blackhole topology. Validation checks the endpoint radio evidence and complete Cooja execution.

\section{Sybil Identity Manipulation}

Sybil is implemented as the new attack surface. After activation, the malicious mote emits RPL DIO messages whose apparent IPv6 source identity rotates across a set of virtual identities. The underlying mote, process, and RPL state remain real. Neighbours therefore observe multiple identities associated with one physical sender.

The design is intentionally bounded. It demonstrates identity manipulation at the RPL control-plane observation surface, but it does not claim to break a cryptographic identity system. Five attack runs produced 150 to 151 spoofed DIO identity events each, while five matched controls produced no activation or spoofing events. The generic IDS does not consume those explicit marker counts. Instead, it uses observable control-message, sender-diversity, radio, and routing-state features.

\section{Feature and Dataset Pipeline}

The routing feature extractor parses Cooja timestamps, application counters, radio records, DIO ranks, parent switches, neighbour state, and topology links. It builds one row for each 60-second interval, including intervals with no event so that missing activity is represented as zero rather than an absent row. Rank summaries include minimum, mean, maximum, and range. Stateful low-rank features carry the latest observed rank by receiver-sender pair across windows, allowing persistent sinkhole evidence to be represented.

The dataset builder reads the already validated feature extracts for blackhole, sinkhole, the six additional attacks, and Sybil. It projects them onto a common schema, adds source-file provenance, and verifies exactly 90 family-mode-seed groups. The generated manifest records each group, run ID, source file, nine window starts, and the 240-second activation time. The audit independently rechecks all constraints and writes a feature-eligibility table.

\section{Static IDS and Adaptation Implementation}

The Gaussian model estimates class means, variances, and priors for the selected features. The CART implementation recursively selects splits that reduce impurity. Both models run locally and require no cloud service. Cross-attack scripts enumerate ordered train and test families, remove metadata and attack markers, train on the source distribution, and calculate confusion-matrix metrics on the target distribution.

The adaptation script enumerates complete target seed combinations. For one target seed there are five possible adaptation choices. For two and three seeds there are ten combinations. Each choice is evaluated only against remaining target seeds. The result table therefore reports both the adaptation size and number of splits, rather than presenting one convenient seed selection.

\section{Online Monitoring}

The CUSUM implementation reduces persistent low-rank state to a scalar score built from non-root low-rank sender, pair, and receiver-exposure counts. Calibration uses source windows that exclude the target seed. Accumulated positive deviation is compared with a fixed threshold. The detector consumes no family label, attack marker, or future window.

DDM monitors the online classification error stream. It tracks the minimum observed error rate and standard deviation, then signals warning or drift when the current estimate exceeds standard thresholds. In this project, errors are delayed by one window to reflect that labels are not necessarily available at prediction time. This explains the expected timing difference between routing-state CUSUM and DDM.

\section{Rank-Parent Defence Prototype}

The defence hook is disabled by default. When enabled, it identifies a candidate as suspicious if the candidate advertises the RPL root rank but its link-local identity does not match the DODAG root. Parent selection first attempts to retain an existing preferred safe parent, then searches for a non-suspicious alternative. Logging is rate limited, but selection decisions can still occur frequently.

Several iterations were preserved rather than erased. Strict rejection and alternative-parent variants generated diagnostic runs. The final stable-parent variant reduced some immediate rejection behaviour but still produced extensive churn. Retaining these failures supports a transparent engineering account: the defect was not a build failure, but an unsafe control rule derived from evidence that was more suitable for detection.

\section{Reproducibility and Portability}

Dataset and analysis scripts use local Python and standard data files. The Cooja applications build against Contiki-NG on the user's macOS environment. No OpenAI key or AI service is required for attack execution, feature extraction, model evaluation, drift detection, or figure generation. The dissertation figures are generated from preserved CSV results and copied into the Overleaf package as SVG files.

\section{Cooja Configuration Generation}

Configuration generators create one Cooja file per family, mode, and seed. They set mote firmware, positions, random seed, radio medium, log plugins, and timeout. Generated names encode these variables, for example \texttt{BH\_ATTACK\_N16\_SEED123456}. This naming convention allows validators and feature extraction scripts to recover family, mode, node count, and seed without maintaining a separate hidden mapping.

Attack and control configurations are matched at the topology and seed level. The malicious application is the principal difference. For the robustness campaign, the condition prefix is added to the run name. The extractor recognises \texttt{ATTACKER\_RELOCATED} and \texttt{LOSSY\_RADIO} while retaining the underlying family and mode.

The Cooja Test Script Editor runs the simulation until the configured timeout and writes \texttt{COOJA.testlog}. Radio events are stored separately in \texttt{COOJA.radio}. A test completion marker distinguishes a successful 540-second simulation from a process that exited early. The final timeout is 541 seconds to avoid racing the last expected log event.

\section{Routing Telemetry Instrumentation}

Generic RPL INFO logs provide received and transmitted DIO ranks, DIS and DAO activity, parent switches, neighbour state, and topology links. This instrumentation is common across attack and control runs and can therefore serve as model evidence. Attack-specific logs are separate and classified as validation markers.

The distinction became important during sinkhole development. A weak flag in an early copied application produced only one false-rank event before normal advertisement resumed. The final application uses the intended strong variable definition, which produces persistent minimum-rank evidence. The invalid preliminary runs remain under a diagnostic directory and are not included in the final feature dataset.

\section{Derived Topology Features}

The extractor maintains the most recent observed parent for each child and derives edge count, number of unique parents, maximum parent fan-out, direct root children, mean depth, and maximum depth. It also stores the most recent received rank for each receiver-sender pair. From this state it derives low-rank pair count, distinct low-rank senders, and the number of receivers exposed to a suspicious low-rank non-root sender.

These features are approximations reconstructed from logs, not a perfect global network state. Missing observations, collisions, and logging level can affect them. Their value lies in exposing persistent routing behaviour that a traffic-only feature set cannot represent.

\section{Result and Figure Generation}

Every primary figure has a matching CSV source. The cross-attack heatmap is generated from CART rows in the nine-family matrix. The adaptation figures use summary rows for zero through three target seeds. Online detection uses the detector comparison table. Robustness uses zero and three target-condition seed summaries. The defence figure uses post-activation means across five matched seeds.

The figure scripts use only the preserved results and do not refit models. This separation reduces the risk that editing a caption or visual design silently changes an experiment. The SVG files are included directly in Overleaf, preserving vector text and lines.

\section{Software Boundaries}

The project uses three different kinds of code. Attack code changes simulated node behaviour. Instrumentation and extraction code observes runs and creates features. Analysis code trains models, evaluates drift and adaptation, and generates tables. Keeping these boundaries visible is important because an attack marker from the first layer must not become a predictive shortcut in the third.

The external Gope scripts form a fourth boundary. They read the supplied CSV files and reproduce static results using that schema. They do not write into the Cooja feature corpus. The offline trust package similarly creates derived interpretation fields without changing the frozen baseline dataset.
